\section{Colorings}

\begin{boxdefinition}[Coloring]
    Given a graph $G = (V, E)$, a $k$-coloring of $G$ is a map $c : V \to [k]$ such that $u \sim v \implies c(u) \ne c(v)$.
\end{boxdefinition}

\subsection{$2$-Colorings of Graphs}

Let's start with a motivating question: which graphs can be 2-colored?

Any odd cycle graph can't be 2-colored. That is, any graph that looks like
\begin{figure}[H]
    \centering
    \begin{tikzpicture}
        \cyclegraph{3}
        \begin{scope}[xshift = 3cm]
            \cyclegraph{5}
        \end{scope}
        \begin{scope}[xshift = 6cm]
            \cyclegraph{7}
        \end{scope}
        \node at (8.25, 0.1) {$\cdots$};
    \end{tikzpicture}
\end{figure}

Let's study when odd cycles exist.

\begin{boxlemma}\label{Ch1:Lemma:Odd_Closed_Walk}
    If $G$ has an odd closed walk, then $G$ has an odd cycle.
\end{boxlemma}
Recall that a closed walk need not be a cycle: walks can contain repetitions. So even line graphs can have closed walks (start from one vertex, go to its neighbor, return).
\begin{proof}
    Suppose $G$ has an odd closed walk. Let $W$ be such a walk of minimal length. Write $W = \parenth{x_1, \ldots, x_i, \ldots, x_j, \ldots, x_{\ell}}$, with $x_{\ell} = x_1$. Suppose $W$ is not a cycle. Then there must be some $x_i, x_j$ such that $x_i = x_j$ even though $i \neq j$, ie, there must be some repetition. Then, $\parenth{x_1, \ldots, x_i, x_{j+1}, \ldots, x_{\ell}}$ and $\parenth{x_i, x_{i+1}, \ldots, x_{j}}$ are both shorter walks than $W$. Moreover, one of them must be of odd length. This contradicts minimality of the length of $W$. So $W$ must be a cycle.
\end{proof}

Here's a cool result.

\begin{boxtheorem}
    A graph $G = (V, E)$ can be $2$-colored if and only if $G$ has no odd cycle.
\end{boxtheorem}
\begin{proof} \hfill
    \begin{description}
        \item[$\parenth{\implies}$] We've already seen that odd cycles can't be $2$-colored. So if $G$ can be $2$-colored, it can't have an odd cycle. We've essentially seen this from the examples above.

        \item[$\parenth{\impliedby}$] We may assume $G$ is connected, since otherwise we can color each component separately. Fix some $v_0 \in V$. Define $c : V \to \set{1, 2}$ by
        \begin{align*}
            c(u) =
            \begin{cases}
                2 & \text{ if } \dist{v_0, u} \text{ is odd} \\
                1 & \text{ if } \dist{v_0, u} \text{ is even}
            \end{cases}
        \end{align*}
        where $\dist{\cdot, \cdot}$ denotes the shortest number of edges from one vertex to another.

        We claim that this $c$ is a valid $2$-coloring.

        Indeed, suppose it is not. Then, there must be some $u, v \in V$ such that $u \sim v$ and $c(u) = c(v)$. Now let's consider the two possible cases on the color.

        Suppose $c(u) = c(v) = 1$ for adjacent $u, v$. Then, there is a path of even length between $v$ and $v_0$ and another path of even length between $u$ and $v_0$. Concatenating the path from $v$ to $v_0$ with the one from $v_0$ to $u$, and closing up with the edge $uv$, we get a closed walk of length even $+$ even $+ 1$. This is odd, so by \Cref{Ch1:Lemma:Odd_Closed_Walk}, $G$ has an odd cycle.

        The case $c(u) = c(v) = 2$ is identical: the two paths are now of odd length, so the closed walk they form with the edge $uv$ has length odd $+$ odd $+ 1$, which is odd again.

        Either way, $G$ has an odd cycle, contrary to hypothesis. So $c$ must be a valid $2$-coloring.
    \end{description}
\end{proof}

\subsection{Colorings of Hypergraphs}

Coloring is not really a question about graphs. All the definition above asks of $G$ is that we know which pairs of vertices may not share a color, and nothing stops us forbidding larger sets instead.

\begin{boxdefinition}[Hypergraph]
    Given a vertex set $V$, a \textbf{hypergraph} is a collection $H$ of subsets of $V$. The sets in $H$ are called \textbf{hyperedges}.
\end{boxdefinition}

\begin{boxdefinition}[Uniformity]
    A hypergraph $H$ is \textbf{$k$-uniform} if $\abs{f} = k$ for all $f \in H$.
\end{boxdefinition}

The following is thus obvious.

\begin{boxproposition}
    A graph is a $2$-uniform hypergraph.
\end{boxproposition}

\begin{boxdefinition}
A proper coloring of a hypergraph $H$ is an assignment $c : V \to [r]$ such that $\forall f \in H$, $\abs{c(f)} \geq 2$.
\end{boxdefinition}

We say ``no edge is monochromatic''.

We can ask ourselves the following question: \textbf{``Which $3$-uniform hypergraphs are $2$-colorable?}

For each $k$, define $m(k)$ to be the smallest $m$ such that $\exists$ some $k$-uniform hypergraph with $m$ edges that is not $2$-colorable. We can show that $m(2) = 3$. But what is $m(3)$? That is, how many $3$-edges do I need to prevent $2$-colorability?

Let $H$ be a $3$-uniform hypergraph that's not $2$-colorable. Assume, further, that $H$ has at most $6$ edges (each being a subset of size $3$).

First, we'll show that we can reduce, without loss of generality, to the case where $H$ has exactly $6$ vertices. Indeed, if $H$ has vertices $x, y$ that are not in any common edge, then I can define $H'$ by ``merging'' $x$ and $y$. So $H'$ has one fewer vertex than $H$. This $H'$ cannot be $2$-colorable: if it were, then that coloring would extend to $H$ by using the color of the merged vertex at both $x$ and $y$, which would make $H$ $2$-colorable. And when $H$ has more than $6$ vertices, such a pair always exists.

An edge has $3$ vertices, so it covers the $\binom{3}{2} = 3$ pairs among them, and $H$ has at most $6$ edges, so at most
\begin{align*}
    6 \cdot \binom{3}{2} = 18
\end{align*}
pairs are covered in total. But if $H$ has $n \geq 7$ vertices, there are $\binom{n}{2} \geq \binom{7}{2} = 21$ pairs to cover. So some pair is left uncovered, and merging it brings us one vertex closer to $6$.

So now assume that $H$ has exactly $6$ edges and $6$ vertices and is not $2$-colorable. 
There is no loss in the vertex count in the other direction either: a hypergraph with fewer than $6$ vertices can be handed as many extra vertices as it likes, and vertices lying in no edge change nothing about which colorings are proper.

Now count colorings. Among all the ways of coloring $6$ vertices with two colors, consider only the balanced ones, which split $V$ into two sets of three; there are $\frac{1}{2}\binom{6}{3} = 10$ of these. Such a coloring fails exactly when some edge is monochromatic, and an edge of $H$ has three vertices, so it is monochromatic under a balanced coloring precisely when it is one of the two color classes. Each edge $f$ therefore rules out exactly one of the ten balanced colorings, namely the one splitting $V$ into $f$ and $V \setminus f$. Six edges rule out at most six of the ten, so at least four survive, and any one of them is a proper $2$-coloring of $H$. This contradicts our assumption, so no such $H$ exists and
\begin{align*}
    m(3) \geq 7
\end{align*}

That bound is attained, and the hypergraph attaining it is a familiar object.

\begin{boxexample}[The Fano Plane]
    The \textbf{Fano plane} is the $3$-uniform hypergraph on the seven vertices below, whose seven edges---by convention, its \textbf{lines}---are the three sides of the triangle, the three medians, and the circle through the three edge midpoints.
\end{boxexample}

\begin{figure}[H]
    \centering
    \begin{tikzpicture}
        \fanoplane
    \end{tikzpicture}
\end{figure}

Every vertex lies on exactly three lines, any two vertices lie on exactly one line together, and any two lines meet in exactly one vertex. That last property is what does the work below.

\begin{boxproposition}
    The Fano plane is not $2$-colorable. Hence $m(3) = 7$.
\end{boxproposition}
\begin{proof}
    Of the $\binom{7}{4} = 35$ sets of four vertices, exactly $28$ contain a line: a line together with any one of the four vertices off it gives such a set, which is $7 \cdot 4 = 28$ sets in all, and no two of them coincide because four vertices cannot contain two lines---two lines meet in a vertex, so between them they cover five. The seven sets left over are the complements of the seven lines.

    Now suppose we had a proper $2$-coloring, with color classes $A$ and $B$, and say $\abs{A} \leq 3$. If $\abs{A} = 3$, then $A$ is not a line, since $A$ is monochromatic; so $B = V \setminus A$ is a set of four vertices that is not the complement of a line, and hence contains one. If $\abs{A} \leq 2$, then the vertices outside $B$ lie on at most $3 + 3 - 1 = 5$ lines between them, so at least two lines avoid them both and lie inside $B$. Either way some line is monochromatic, which is a contradiction.

    So the Fano plane is a $3$-uniform hypergraph with seven edges that is not $2$-colorable, giving $m(3) \leq 7$.
\end{proof}
